\paragraph{RQ1：构建与容量。}
如何把JHQ层次训练、编码并组织成可搜索的GPU常驻索引，实际构建成本和容量是多少？
证据为物理布局、train/add记录和常驻核算（\cref{sec:7-7}）。
构建失败或分配超过设备预算会限制容量主张，常驻样本不代表构建峰值。

\paragraph{RQ2：正确高效的执行。}
哪些面向表示的评分、保留和调度选择，在满足分数与并列契约时改善GPU搜索？
LUT、packed load、精排和扫描配对实验连接参考检查与batch敏感性
（\cref{sec:7-3,sec:7-4,sec:7-6}）。非并列选择错误挑战正确性，持续减速挑战执行选择。

\paragraph{RQ3：预算质量与总成本。}
无标签一致性能否选出有用预算，什么抽样与复用条件下校准值得？
互斥查询评价固定预算、抽样／FULL、枚举及GT oracle（\cref{sec:7-5}）。
越阈挑战可靠性，包含校准后没有正节省则挑战净收益。
oracle是评估对照，不是可部署的无标签决策。

\paragraph{RQ4：竞争区间。}
完整系统在共同主机到主机边界下，哪些实测recall区间具有竞争力？
依据独立调优E1前沿和明确显存定义（\cref{sec:7-2}）。
int8图对手、高recall反转和未覆盖区间都是答案的一部分，
不追求跨硬件、显存预算和请求大小的普遍排名。

第3章将这些问题转成受控比较，分别回答建成了什么、执行是否正确、
工作负载需要多少工作，以及完整系统有没有有用竞争区间。
